The primary threat to external validity resides in the selection criteria and
the representativeness of the collected issues.
%
We are certain that all 30 issues are indeed soundness issues because we studied
each in depth;
%
readers can also validate this claim by reviewing the code snippets in our
supplementary material~\cite{supp} or the issues themselves on GitHub, as the
dataset size is manageable.
%
However, additional soundness issues may exist that were missed due to our
conservative selection criteria.
%
The automatic filtering process relied on GitHub labels;
%
although these labels are routinely assigned by \rustc{} developers, some issues
might not have been properly labeled at the time of collection.
%
While we followed the classification used in prior work~\cite{liu2025empiricala}
to identify labels related to type checking, other relevant labels might have
been overlooked.
%
Moreover, we filtered out issues with labels that appeared irrelevant to
soundness, but it is possible that a soundness issue also revealed other
problems and was thus excluded.
%
Notably, as mentioned in \cref{sec:method}, 7 issues in the dataset of the
previous work were excluded by our filtering due to their labels, confirming the
risk of missing soundness issues.
%
Even though we complemented our dataset with these 7 specific issues, other
soundness issues might have been missed for the same reason.
%
Finally, our manual scan retained issues only if they provided clear evidence
that the expected behavior was rejection, which could exclude less-documented
cases.
%
Nevertheless, even if the dataset is not exhaustive, we believe our findings
remain valid as they are derived from actual soundness issues.

Regarding this threat, we conducted a lightweight study of issues excluded due
to the I-ICE label, which indicates that the compiler crashes because of an
internal compiler error (ICE).
%
We found that 181 excluded issues have the I-ICE label, while only one of them
also has the I-unsound label.
%
As a soundness issue is not necessarily marked with I-unsound, this does not
strictly mean that all except one are not soundness issues.
%
That said, we believe this result suggests that most I-ICE issues are unlikely
to be soundness issues.
%
The primary reason is that an ICE prevents code from compiling, while soundness
issues require code that should not type-check to compile successfully.

Another threat to external validity stems from the exclusion of issues related
to feature-gated features.
%
Inspecting such issues would reveal additional soundness issues, especially
considering that unstable features likely have more compiler bugs than stable
features.
%
We found that among 223 issues excluded due to an ``F-'' label, 10 issues also
have I-unsound, corroborating this expectation.
%
However, we believe that focusing on soundness issues within stable features
helps better understand the severity of soundness issues in \rustc{}.

A threat to construct validity concerns the accuracy of our characterization of
features, symptoms, consequences, and bug status.
%
The identification of these attributes is inherently susceptible to researcher
bias or human error.
%
To mitigate this, we employed three strategies.
%
First, two authors independently investigated the issues and conducted extensive
discussions to reach a consensus on the characterization of each issue.
%
Second, we cross-referenced each issue with multiple external resources,
including related issues, PRs, and discussions on the official Rust Zulip
channels, as well as documents such as the Reference, the FLS, and RFCs.
%
Third, we designed objective criteria for bug status and checked the status by
type-checking each code using the latest \rustc{} release and inspecting issue
labels and their assigners.

Finally, a threat to conclusion validity lies in the small size of the dataset,
which limits the generalizability of our findings.
%
Nevertheless, we believe our work provides value to the community.
%
Our main contribution is reporting that \rustc{} has soundness bugs despite
Rust's emphasis on type soundness, highlighting the urgent need for research
community efforts.
%
Our in-depth, multidimensional analysis yields actionable insights that
facilitate future research on testing the Rust type checker and formalizing the
type system.
